\documentclass[12pt,a4paper]{report}
\usepackage[utf8]{inputenc}
\usepackage[T2A]{fontenc}
\usepackage[russian]{babel}
\usepackage{amsmath,amssymb,amsthm,mathtools}
\usepackage{geometry}
\geometry{top=20mm,bottom=20mm,left=25mm,right=15mm}
\usepackage{hyperref}
\hypersetup{
colorlinks=true, linkcolor=blue!70!black, citecolor=green!50!black, urlcolor=blue!80!black }
\usepackage{booktabs}
\usepackage{array}
\usepackage{xcolor}
\usepackage{amsopn}
\usepackage{amsfonts}
%\usepackage{unicode-math}
\usepackage{amstext}

\theoremstyle{definition}
\newtheorem{question}{Question}
\newtheorem{solution}{Solution}

\theoremstyle{plain}
\newtheorem{theorem}{Theorem}[chapter]
\newtheorem{lemma}[theorem]{Lemma}

\DeclareMathOperator{\spanop}{span}
\DeclareMathOperator{\rank}{rank}
\DeclareMathOperator{\nullity}{nullity}
\DeclareMathOperator{\im}{im}
\DeclareMathOperator{\diag}{diag}
\DeclareMathOperator{\tr}{tr}
\DeclareMathOperator{\Col}{Col}
\DeclareMathOperator{\Row}{Row}
\DeclareMathOperator{\Null}{Null}

\title{\textbf{\LARGE Полные ответы на экзаменационные билеты}\\[4mm]
\Large Подготовка к экзамену по курсу «Refresher: Линейная алгебра»}
\author{\textbf{Ответы на все теоретические, практические и базовые вопросы}}
\date{2026}

\begin{document}

    \maketitle

    \tableofcontents
    \newpage

    \chapter*{Структура экзамена и методические указания}
    \addcontentsline{toc}{chapter}{Структура экзамена и методические указания}

Экзаменационные материалы курса структурированы по трем категориям сложности и четырем тематическим модулям:
    \begin{enumerate}
        \item \textbf{Раздел Theory (Теоретические вопросы):} Требуют глубокого концептуального понимания, вывода
        фундаментальных формул и строгих математических доказательств
        (индукция, от противного, спектральная декомпозиция).
        \item \textbf{Раздел Practice (Практические вопросы):} Направлены на вычислительные методы, пошаговые алгоритмы,
        практический анализ матричных разложений, диагностику размерностей и прикладные модели (PCA, рекомендательные
        системы, решение СЛУ).
        \item \textbf{Раздел Basic (Базовые вопросы):} Охватывают канонические определения, базовые аксиомы,
        формулировки ключевых теорем и минимальный математический аппарат, необходимый для сдачи курса.
    \end{enumerate}

В данном пособии приведены исчерпывающие, строгие и наглядные ответы на каждый из 57 экзаменационных вопросов.

    \newpage

% =========================================================================
% ЧАСТЬ 1: ТЕОРИЯ (THEORY)
% =========================================================================

    \chapter{Теоретические вопросы (Theory)}
\section{Модуль I. Векторные пространства и евклидова геометрия}
% -------------------------------------------------------------------------

\subsection{Теоретический вопрос I.1: Почему координаты в базисе единственны?}
\textbf{Формулировка вопроса:} Докажите, почему координаты любого вектора в фиксированном базисе векторного пространства
определены однозначно.

\textbf{Ответ и строгое доказательство:}
Пусть $V$ — вещественное векторное пространство, и $\mathcal{B} = (e_1, e_2, \dots, e_n)$ — базис в $V$.
По определению базиса выполнены два фундаментальных условия:
\begin{enumerate}
    \item Система векторов $(e_1, \dots, e_n)$ \textbf{порождает} все пространство: $\spanop(e_1, \dots, e_n) = V$.
    \item Система векторов $(e_1, \dots, e_n)$ \textbf{линейно независима}: равенство $\sum_{i=1}^n c_i e_i = 0$
    возможно тогда и только тогда, когда $c_1 = c_2 = \dots = c_n = 0$.
\end{enumerate}

Возьмём произвольный вектор $x \in V$.
Из первого свойства следует, что $x$ представим в виде линейной комбинации базисных векторов:
\[ x = x_1 e_1 + x_2 e_2 + \dots + x_n e_n = \sum_{i=1}^n x_i e_i, \quad x_i \in \mathbb{R}. \]
Числа $(x_1, \dots, x_n)$ называются координатами вектора $x$ в базисе $\mathcal{B}$.

\textbf{Доказательство единственности (от противного):} Предположим, что для того же самого вектора $x$ существует
второй набор координат $(y_1, y_2, \dots, y_n)$, отличный от первого:
\[ x = y_1 e_1 + y_2 e_2 + \dots + y_n e_n = \sum_{i=1}^n y_i e_i. \]
Вычтем второе векторное равенство из первого:
\[ 0 = x - x = \sum_{i=1}^n x_i e_i - \sum_{i=1}^n y_i e_i = \sum_{i=1}^n (x_i - y_i) e_i. \]
Мы получили линейную комбинацию базисных векторов $e_1, \dots, e_n$, равную нулевому вектору $0$, с коэффициентами $c_i
= (x_i - y_i)$.
Но в силу второго свойства базиса векторы $e_1, \dots, e_n$ \textbf{линейно независимы}.
Следовательно, эта линейная комбинация обязана быть тривиальной, то есть все коэффициенты равны нулю:
\[ x_i - y_i = 0 \iff x_i = y_i \quad \forall i \in \{1, 2, \dots, n\}. \]
Таким образом, $x_1 = y_1, x_2 = y_2, \dots, x_n = y_n$, что доказывает абсолютную единственность координат.

\textbf{Концептуальный вывод:} Единственность координат гарантируется свойством \emph{линейной независимости}.
Если бы базисные векторы были зависимы, существовало бы ненулевое ядро отображения координат, что привело бы к
неединственности разложения.

% -------------------------------------------------------------------------

\subsection{Теоретический вопрос I.2: Выведите формулу проекции из ортогональности}
\textbf{Формулировка вопроса:} Выведите формулу ортогональной проекции вектора на вектор и обобщите ее на случай
проекции на подпространство, опираясь исключительно на геометрическое условие ортогональности ошибки.

\textbf{Ответ и строгий вывод:}

\textbf{1.
Проекция вектора на вектор (одномерный случай):} Пусть даны ненулевой вектор $u \in V$ ($u \ne 0$) и вектор $v \in V$.
Мы ищем ортогональную проекцию вектора $v$ на прямую $\spanop(u)$, то есть вектор $p = P_u(v)$.
\begin{enumerate}
    \item Так как $p$ лежит на прямой, натянутой на $u$, он коллинеарен $u$:
    \[ p = \alpha u \quad \text{для некоторого неизвестного скаляра } \alpha \in \mathbb{R}. \]
    \item Определим вектор невязки (ошибки проекции):
    \[ e \coloneqq v - p = v - \alpha u. \]
    \item По геометрическому смыслу ортогональной проекции вектор ошибки $e$ должен быть строго \textbf{ортогонален}
    направлению проецирования $u$:
    \[ e \perp u \iff \langle e, u \rangle = 0 \iff \langle v - \alpha u, u \rangle = 0. \]
    \item Используя билинейность скалярного произведения:
    \[ \langle v, u \rangle - \alpha \langle u, u \rangle = 0 \iff \alpha \langle u, u \rangle = \langle v,
    u \rangle.\]
    \item Так как $u \ne 0$, квадрат его нормы $\langle u, u \rangle = \|u\|^2 > 0$.
    Отсюда однозначно находим скаляр $\alpha$:
    \[ \alpha = \frac{\langle v, u \rangle}{\langle u, u \rangle} = \frac{\langle v, u \rangle}{\|u\|^2}. \]
    \item Подставляя $\alpha$ в выражение для $p$, получаем каноническую \textbf{формулу ортогональной проекции}:
    \[ P_u(v) = \frac{\langle v, u \rangle}{\langle u, u \rangle} u = \frac{\langle v, u \rangle}{\|u\|^2} u. \]
\end{enumerate}

\textbf{2.
Обобщение: Проекция на подпространство (матричный вывод):} Пусть $U \le \mathbb{R}^n$ — подпространство размерности $k$,
натянутое на столбцы матрицы $A \in \mathbb{R}^{n \times k}$ полного ранга ($\rank A = k$).
Проекция вектора $y \in \mathbb{R}^n$ на $U = \Col(A)$ представляется в виде $p = A x$ для некоторого $x \in
\mathbb{R}^k$.
Вектор ошибки равен $e = y - p = y - A x$.
Условие ортогональности проекции: вектор $e$ должен быть ортогонален каждому столбцу матрицы $A$, то есть:
\[ A^\top e = 0 \iff A^\top (y - A x) = 0 \iff A^\top A x = A^\top y. \]
Поскольку $\rank A = k$, квадратная матрица Грама $A^\top A \in \mathbb{R}^{k \times k}$ строго положительно определена
и обратима.
Находим $x$:
\[ x = (A^\top A)^{-1} A^\top y. \]
Следовательно, вектор проекции равен:
\[ p = A x = A (A^\top A)^{-1} A^\top y. \]
Отсюда получаем матрицу ортогонального проектора $P = A (A^\top A)^{-1} A^\top$, для которой автоматически выполнены
свойства идемпотентности ($P^2 = P$) и симметричности ($P^\top = P$).

% -------------------------------------------------------------------------

\subsection{Теоретический вопрос I.3:Объясните иерархию metric $\to$ norm $\to$ inner product $\to$ Euclidean space}
\textbf{Формулировка вопроса:} Объясните строгую математическую иерархию геометрических структур от метрических
пространств до евклидова пространства.
Приведите условия перехода и контрпримеры.

\textbf{Ответ:}

Иерархия строится по принципу последовательного усиления аксиоматических требований:
\[ \text{Метрическое пространство} \supset \text{Нормированное векторное пространство} \supset
\text{Предгильбертово пространство} \supset \text{Евклидово пространство}.
\]

\textbf{1.
Метрическое пространство $(M, d)$:} Множество $M$ с функцией расстояния $d(x, y)$, удовлетворяющей аксиомам:
\begin{itemize}
    \item $d(x, y) \ge 0$, причём $d(x, y) = 0 \iff x = y$;
    \item $d(x, y) = d(y, x)$ (симметрия);
    \item $d(x, z) \le d(x, y) + d(y, z)$ (неравенство треугольника).
\end{itemize}
Здесь элементы не обязательно складывать и умножать на числа — есть только расстояние.

\textbf{2.
Переход «Метрика $\to$ Норма»:} Если на множестве введена структура векторного пространства $V$ над $\mathbb{R}$, можно
ввести \textbf{норму} $\|x\|$ (длину вектора):
\begin{itemize}
    \item $\|x\| \ge 0$, причём $\|x\| = 0 \iff x = 0$;
    \item $\|\alpha x\| = |\alpha| \|x\|$ (\textbf{абсолютная однородность});
    \item $\|x + y\| \le \|x\| + \|y\|$ (\textbf{субаддитивность}).
\end{itemize}
\emph{Индукция метрики:} Всякая норма порождает метрику расстояния по формуле $d(x, y) \coloneqq \|x - y\|$.
Такая метрика инвариантна к сдвигам: $d(x+z, y+z)=d(x, y)$ и однородна: $d(\alpha x,\alpha y) = |\alpha| d(x, y)$.\\
\emph{Контрпример (почему не всякая метрика есть норма):} Дискретная метрика $d(x, y) = 1$ при $x \ne y$ не
удовлетворяет свойству однородности: $d(2x, 2y) = 1 \ne 2 d(x, y)$.

\textbf{3.
Переход «Норма $\to$ Скалярное произведение»:} Скалярное произведение $\langle x, y \rangle$ — это билинейная,
симметричная и положительно определенная форма.\\ \emph{Индукция нормы:} Всякое скалярное произведение порождает
евклидову норму:
\[ \|x\| \coloneqq \sqrt{\langle x, x \rangle}. \]
Неравенство треугольника для этой нормы гарантируется неравенством Коши-Буняковского-Шварца ($|\langle x, y \rangle| \le
\|x\| \|y\|$).\\ \emph{Критерий обратного перехода (Теорема фон Неймана-Жордана):} Норма порождается скалярным
произведением тогда и только тогда, когда для любых векторов выполнено \textbf{тождество параллелограмма}:
\[ \|x + y\|^2 + \|x - y\|^2 = 2\|x\|^2 + 2\|y\|^2. \]
\emph{Контрпример (почему не всякая норма порождается скалярным произведением):} В $\mathbb{R}^2$ норма $L_1$: $\|x\|_1
= |x_1| + |x_2|$.
Возьмём $x = (1, 0)$ и $y = (0, 1)$.
Тогда $\|x\|_1 = 1$, $\|y\|_1 = 1$.
Имеем $x+y = (1, 1)$ , $\|x+y\|_1 = 2$; $x-y = (1, -1)$, $\|x-y\|_1 = 2$.
Слева: $2^2 + 2^2 = 8$.
Справа: $2(1^2) + 2(1^2) = 4$.
Так как $8 \ne 4$, тождество параллелограмма нарушено.
Следовательно, в пространстве с нормой $L_1$ понятия угла и скалярного произведения не существует.

\textbf{4.
Евклидово пространство:} Это \textbf{конечномерное} вещественное векторное пространство со скалярным произведением.
В нём одновременно существуют все четыре согласованных понятия: скалярное произведение $\langle x, y \rangle$, евклидова
норма $\|x\| = \sqrt{\langle x, x \rangle}$, евклидово расстояние $d(x, y) = \|x - y\|$ и угол $\cos \theta =
\frac{\langle x, y \rangle}{\|x\| \|y\|}$.

\newpage

\section{Модуль II. Матрицы, линейные отображения, ранг и СЛУ}
% -------------------------------------------------------------------------

\subsection{Теоретический вопрос II.1: Почему columns of a matrix являются images of basis vectors?}
\textbf{Формулировка вопроса:} Докажите, почему столбцы матрицы линейного отображения в точности равны координатным
представлениям образов базисных векторов.

\textbf{Ответ и строгое доказательство:}

Пусть $T \colon V \to W$ — линейное отображение между векторными пространствами.
Зафиксируем базис $\mathcal{B} = (v_1, \dots, v_n)$ в $V$ и базис $\mathcal{C} = (w_1, \dots, w_m)$ в $W$.
Возьмём вектор $x \in V$.
По теореме о базисе он единственным образом представляется в виде:
\[ x = \sum_{j=1}^n x_j v_j, \quad \text{где } [x]_\mathcal{B} = \begin{pmatrix}
                                                                     x_1 \\ \vdots \\ x_n
\end{pmatrix}. \]
Применим отображение $T$ к вектору $x$.
В силу свойства \textbf{линейности} оператора $T$:
\[ T(x) = T\left(\sum_{j=1}^n x_j v_j\right) = \sum_{j=1}^n x_j T(v_j). \]
Каждый вектор $T(v_j) \in W$ раскладывается по базису $\mathcal{C}$:
\[ T(v_j) = \sum_{i=1}^m a_{ij} w_i, \quad \text{то есть } [T(v_j)]_\mathcal{C} =
\begin{pmatrix}
            a_{1j} \\ a_{2j} \\ \vdots \\ a_{mj}
\end{pmatrix} \].
Подставим это выражение в формулу для $T(x)$:
\[ T(x) = \sum_{j=1}^n x_j \left( \sum_{i=1}^m a_{ij} w_i \right) = \sum_{i=1}^m \left( \sum_{j=1}^n a_{ij} x_j
\right) w_i.
\] Следовательно , $i$-я координата вектора $T(x)$ в базисе $\mathcal{C}$ равна $y_i = \sum_{j=1}^n a_{ij} x_j$.
В терминах матричного произведения это тождественно равно:
\[ \begin{pmatrix}
       y_1 \\ \vdots \\ y_m
\end{pmatrix} = \begin{pmatrix}
                    a_{11} & a_{12} & \dots & a_{1n} \\
                    a_{21} & a_{22} & \dots & a_{2n} \\
                    \vdots & \vdots & \ddots & \vdots \\
                    a_{m1} & a_{m2} & \dots & a_{mn}
\end{pmatrix} \begin{pmatrix}
                  x_1 \\ \vdots \\ x_n
\end{pmatrix} = A [x]_\mathcal{B}. \]

\textbf{Вывод:} Если мы возьмём в качестве $x$ $j$-й базисный вектор $v_j$, его координатный столбец $[v_j]_\mathcal{B}
= e_j = (0, \dots, 1, \dots, 0)^\top$ содержит единицу на $j$-м месте и нули на остальных.
Тогда:
\[ [T(v_j)]_\mathcal{C} = A e_j = A_{\cdot j} \quad (\text{$j$-й столбец матрицы } A). \]
Таким образом , $j$-й столбец матрицы отображения обязан в точности совпадать с вектором координат образа $j$-го
базисного вектора.

% -------------------------------------------------------------------------

\subsection{Теоретический вопрос II.2: Концептуально объясните rank-nullity}
\textbf{Формулировка вопроса:} Объясните концептуальный смысл теоремы о ранге и дефекте ($\dim V = \rank T + \nullity
T$) как закона сохранения размерности.

\textbf{Ответ:}

Теорема о ранге и дефекте формулируется следующим образом:
\[ \dim V = \dim(\ker T) + \dim(\im T) \iff \dim V = \nullity(T) + \rank(T). \]
Для матрицы $A \in \mathbb{R}^{m \times n}$ с $n$ столбцами:
\[ n = \dim(\Null A) + \rank(A). \]

\textbf{Концептуальная интерпретация («Закон сохранения степеней свободы»):}
\begin{enumerate}
    \item Пространство прообразов $V$ обладает $\dim V = n$ независимыми степенями свободы (направлениями движения).
    \item При действии любого линейного оператора $T$ эти степени свободы расщепляются на две строго
    непересекающиеся группы:
    \begin{itemize}
        \item Направления, которые отображение \textbf{уничтожает (схлопывает в нуль)}.
        Это ядро оператора $\ker T$.
        Количество схлопнутых измерений равно дефекту $\nullity(T)$.
        \item Направления, которые \textbf{сохраняются (переживают трансформацию)}.
        Они разворачиваются в образе $\im T$, формируя его размерность — ранг $\rank(T)$.
    \end{itemize}
    \item \textbf{Закон сохранения:} Измерение не может возникнуть из ниоткуда и не может исчезнуть бесследно:
    \[ \text{Исходные измерения} = \text{Потерянные измерения} + \text{Выжившие измерения}. \]
\end{enumerate}

\textbf{Экстремальные случаи:}
\begin{itemize}
    \item Если $\nullity(T) = 0$ (ядро тривиально, $\ker T = \{0\}$), то оператор инъективен: ни одно измерение
    не потеряно, $\rank(T) = n$.
    \item Если $\rank(T) = 0$ (нулевой оператор), все $n$ измерений схлопываются в точку $0$, $\nullity(T) = n$.
\end{itemize}

% -------------------------------------------------------------------------

\subsection{Теоретический вопрос II.3: Почему elementary row operations не искажают решения системы уравнений?}
\textbf{Формулировка вопроса:} Докажите, почему элементарные строчные преобразования строго сохраняют множество решений
системы линейных уравнений.

\textbf{Ответ и строгое доказательство:}

Рассмотрим систему линейных уравнений в матричной форме:
\[ Ax = b, \quad A \in \mathbb{R}^{m \times n}, \; x \in \mathbb{R}^n, \; b \in \mathbb{R}^m. \]
Обозначим множество ее решений через $S = \{ x \in \mathbb{R}^n \mid Ax = b \}$.
Элементарные преобразования строк включают: 1) перестановку строк $i$ и $j$; 2) умножение строки на $c \ne 0$; 3)
прибавление к строке $i$ строки $j$, умноженной на $c$.

\textbf{Лемма об элементарных матрицах:} Каждое элементарное строчное преобразование расширенной матрицы $(A \mid b)$ в
точности эквивалентно умножению ее слева на соответствующую \textbf{элементарную матрицу} $E \in \mathbb{R}^{m \times
m}$:
\[ E (A \mid b) = (EA \mid Eb). \]
Преобразованная система имеет вид:
\[ (EA) x = Eb. \]
Обозначим множество ее решений через $S' = \{ x \in \mathbb{R}^n \mid (EA)x = Eb \}$.

\textbf{Ключевое свойство:} Любая элементарная матрица $E$ \textbf{строго обратима}, причём матрица $E^{-1}$ сама
является элементарной матрицей того же типа:
\begin{itemize}
    \item Для перестановки строк $E_{ij}^{-1} = E_{ij}$;
    \item Для масштабирования $E_i(c)^{-1} = E_i(1/c)$ при $c \ne 0$;
    \item Для сдвига $E_{ij}(c)^{-1} = E_{ij}(-c)$.
\end{itemize}

\textbf{Доказательство эквивалентности множеств решений($S = S'$):}
\begin{enumerate}
    \item \textbf{Включение $S \subseteq S'$:} Пусть $x \in S$, то есть $Ax = b$.
    \item Умножим обе части равенства слева на матрицу $E$:
    \[ E(Ax) = Eb \iff (EA)x = Eb \implies x \in S'. \]
    \item \textbf{Включение $S' \subseteq S$:} Пусть $x \in S'$, то есть $(EA)x = Eb$.
    В силу существования $E^{-1}$ умножим обе части равенства слева на $E^{-1}$:
    \[ E^{-1} ((EA)x)=E^{-1} (Eb) \iff (E^{-1} E) Ax=(E^{-1} E) b \iff I_m Ax=I_m b \iff Ax = b \implies x \in S. \]
\end{enumerate}
Поскольку $S \subseteq S'$ и $S' \subseteq S$, множества решений тождественно совпадают: $S = S'$.
Ни одно решение не теряется и ни одно постороннее решение не появляется.

\newpage

\section{Модуль III. Собственные значения и матричные разложения}
% -------------------------------------------------------------------------

    \subsection{Теоретический вопрос III.1: Выведите characteristic equation for eigenvalues}
\textbf{Формулировка вопроса:} Выведите характеристическое уравнение для собственных значений матрицы из определения
собственного вектора.

\textbf{Ответ и строгий вывод:}

По определению ненулевой вектор $v \in \mathbb{C}^n$ ($v \ne 0$) называется собственным вектором матрицы $A \in
\mathbb{R}^{n \times n}$, отвечающим собственному значению $\lambda \in \mathbb{C}$, если: \[ A v = \lambda v.
\] Перенесем все слагаемые в левую часть, используя единичную матрицу $I_n$: \[ A v - \lambda I_n v = 0 \iff (A -
\lambda I_n) v = 0.
\] Мы получили однородную систему линейных уравнений относительно вектора $v$ с квадратной матрицей коэффициентов
$M(\lambda) \coloneqq A - \lambda I_n \in \mathbb{C}^{n \times n}$.

Условие существования \textbf{ненулевого} собственного вектора $v \ne 0$: Однородная система $M(\lambda) v = 0$ имеет
нетривиальное решение тогда и только тогда, когда ядро матрицы $M(\lambda)$ содержит ненулевые векторы: \[ \ker(A -
\lambda I_n) \ne \{0\} \iff \nullity(A - \lambda I_n) > 0.
\] По теореме о ранге и дефекте это равносильно тому, что ранг матрицы строго меньше ее размера: \[ \rank(A - \lambda
I_n) < n.
\] А квадратная матрица неполного ранга является \textbf{вырожденной}, что по критерию обратимости эквивалентно
равенству нулю ее определителя: \[ \det(A - \lambda I_n) = 0.
\] Многочлен $p_A(\lambda) \coloneqq \det(A - \lambda I_n)$ является многочленом степени $n$ относительно переменной
$\lambda$ и называется \textbf{характеристическим многочленом}, а уравнение $\det(A - \lambda I_n) = 0$ —
\textbf{характеристическим уравнением}.

% -------------------------------------------------------------------------

    \subsection{Теоретический вопрос III.2: Почему distinct eigenvalues дают independent eigenvectors?}
\textbf{Формулировка вопроса:} Докажите теорему: собственные векторы, отвечающие попарно различным собственным
значениям, линейно независимы.

\textbf{Ответ и строгое доказательство:}

    \begin{theorem}
Пусть $v_1, v_2, \dots, v_k$ — собственные векторы матрицы $A$, отвечающие попарно различным собственным значениям
$\lambda_1, \lambda_2, \dots, \lambda_k$ ($\lambda_i \ne \lambda_j$ при $i \ne j$).
Тогда система векторов $\{v_1, \dots, v_k\}$ линейно независима.
    \end{theorem}
    \begin{proof}
Доказательство проведем методом математической индукции по числу векторов $k$.

\textbf{База индукции ($k = 1$):} По определению собственный вектор ненулевой: $v_1 \ne 0$.
Один ненулевой вектор всегда образует линейно независимую систему ($c_1 v_1 = 0 \implies c_1 = 0$).
База доказана.

\textbf{Индукционный переход:} Пусть утверждение справедливо для $k - 1$ векторов: любые $k - 1$ собственных векторов с
различными собственными значениями линейно независимы.
Рассмотрим систему из $k$ векторов и составим их линейную комбинацию, равную нулю:
        \begin{equation} \label{eq:t32_1} c_1 v_1 + c_2 v_2 + \dots + c_{k-1} v_{k-1} + c_k v_k = 0. \end{equation}
Применим оператор $A$ к обеим частям равенства \eqref{eq:t32_1}.
По линейности и свойству $A v_i=\lambda_i v_i$:
        \begin{equation} \label{eq:t32_2}
            c_1 \lambda_1 v_1 + c_2 \lambda_2 v_2 + \dots + c_{k-1} \lambda_{k-1} v_{k-1} + c_k \lambda_k v_k = 0.
        \end{equation}
Теперь умножим исходное уравнение \eqref{eq:t32_1} на скаляр $\lambda_k$:
        \begin{equation} \label{eq:t32_3}
            c_1 \lambda_k v_1 + c_2 \lambda_k v_2 + \dots + c_{k-1} \lambda_k v_{k-1} + c_k \lambda_k v_k = 0.
        \end{equation}
Вычтем почленно равенство \eqref{eq:t32_3} из равенства \eqref{eq:t32_2}: \[ c_1 (\lambda_1 - \lambda_k) v_1 + c_2
(\lambda_2 - \lambda_k) v_2 + \dots + c_{k-1} (\lambda_{k-1} - \lambda_k) v_{k-1} + 0 = 0.
\] Мы получили равенство нулю линейной комбинации первых $k - 1$ векторов $\{v_1, \dots, v_{k-1}\}$.
По индукционному предположению эта система линейно независима.
Следовательно, все ее коэффициенты обязаны быть равны нулю: \[ c_i (\lambda_i - \lambda_k) = 0 \quad \text{для всех } i
\in \{1, 2, \dots, k-1\}.
\] Так как по условию все собственные значения попарно различны, разность $\lambda_i - \lambda_k \ne 0$ при $i< k$.
Отсюда получаем: \[ c_1 = c_2 = \dots = c_{k-1} = 0.
\] Подставим найденные нулевые коэффициенты обратно в исходное равенство \eqref{eq:t32_1}: \[ 0 + 0 + \dots + 0 + c_k
v_k = 0 \iff c_k v_k = 0.
\] Поскольку собственный вектор $v_k \ne 0$, получаем $c_k = 0$.
Итак, все коэффициенты $c_1 = c_2 = \dots = c_k = 0$.
Система векторов $\{v_1, \dots, v_k\}$ линейно независима, что завершает шаг индукции.
    \end{proof}

% -------------------------------------------------------------------------

    \subsection{Теоретический вопрос III.3: Объясните diagonalization through eigenbasis}
\textbf{Формулировка вопроса:} Объясните процедуру и геометрический смысл диагонализации матрицы через базис из
собственных векторов.

\textbf{Ответ:}

Пусть квадратная матрица $A \in \mathbb{R}^{n \times n}$ обладает $n$ линейно независимыми собственными векторами $v_1,
v_2, \dots, v_n$, отвечающими собственным значениям $\lambda_1, \dots, \lambda_n$: \[ A v_j = \lambda_j v_j, \quad j =
1, \dots, n.
\] Составим матрицу $P \in \mathbb{R}^{n \times n}$, столбцами которой являются эти собственные векторы: \[ P =
\begin{pmatrix} | & | & & | \\ v_1 & v_2 & \dots & v_n \\ | & | & & |
    \end{pmatrix}. \]
Поскольку столбцы матрицы $P$ линейно независимы, матрица $P$ невырождена ($\det P \ne 0$) и существует обратная матрица
$P^{-1}$.

Умножим матрицу $A$ на матрицу $P$ по правилу блочного умножения: \[ A P = A \begin{pmatrix} | & & | \\ v_1 & \dots &
v_n \\ | & & |
    \end{pmatrix} = \begin{pmatrix}
| & & | \\ A v_1 & \dots & A v_n \\ | & & |
    \end{pmatrix} = \begin{pmatrix}
| & & | \\ \lambda_1 v_1 & \dots & \lambda_n v_n \\ | & & |
    \end{pmatrix}. \]
С другой стороны, умножим матрицу $P$ справа на диагональную матрицу $\Lambda = \diag(\lambda_1, \dots, \lambda_n)$: \[
P \Lambda = \begin{pmatrix} | & & | \\ v_1 & \dots & v_n \\ | & & |
    \end{pmatrix} \begin{pmatrix}
\lambda_1 & & 0 \\ & \ddots & \\ 0 & & \lambda_n
    \end{pmatrix} = \begin{pmatrix}
| & & | \\ \lambda_1 v_1 & \dots & \lambda_n v_n \\ | & & |
    \end{pmatrix}. \]
Столбцы обеих матриц тождественно равны, следовательно: \[ A P = P \Lambda.
\] Умножив равенство справа на $P^{-1}$, получаем формулу \textbf{диагонализации}: \[ A = P \Lambda P^{-1}.
\] А умножив слева на $P^{-1}$, получаем формулу перехода к диагональному виду: \[ \Lambda = P^{-1} A P. \]

\textbf{Геометрический смысл диагонализации:} Матрица $P$ является матрицей перехода от стандартного базиса к
собственному базису $\mathcal{B}=(v_1, \dots, v_n)$.
Преобразование $y = Ax$ в собственном базисе распадается на $n$ независимых одномерных масштабирований: \[
[y]_\mathcal{B} = \Lambda [x]_\mathcal{B} \iff \begin{pmatrix} y'_1 \\ \vdots \\ y'_n
    \end{pmatrix} = \begin{pmatrix}
\lambda_1 x'_1 \\ \vdots \\ \lambda_n x'_n
    \end{pmatrix}. \]
Собственные направления — это оси естественной системы координат оператора, вдоль которых отсутствует перемешивание
координат.

    \newpage

\section{Модуль IV. SVD и итоговые инструменты}
% -------------------------------------------------------------------------

\subsection{Теоретический вопрос IV.1: Выведите SVD из произведения матриц $A^\top A$}
\textbf{Формулировка вопроса:} Конструктивно выведите сингулярное разложение (SVD) произвольной вещественной матрицы $A
\in \mathbb{R}^{m \times n}$ через спектральный анализ матрицы $A^\top A$.

\textbf{Ответ и строгий вывод:}

Пусть дана произвольная вещественная матрица $A \in \mathbb{R}^{m \times n}$ ранга $\rank(A) = r \le \min(m, n)$.
Рассмотрим квадратную матрицу $S \coloneqq A^\top A \in \mathbb{R}^{n \times n}$.
\begin{enumerate}
    \item \textbf{Свойства матрицы $A^\top A$:}
    \begin{itemize}
        \item Симметричность: $(A^\top A)^\top = A^\top A$.
        \item Неотрицательная определенность: для любого $x \in \mathbb{R}^n$ имеем $x^\top (A^\top A) x =
        (Ax)^\top (Ax) = \|Ax\|_2^2 \ge 0$.
    \end{itemize}
    \item По вещественной спектральной теореме матрица $A^\top A$ обладает ортонормированным базисом из
    собственных векторов $(v_1, v_2, \dots, v_n)$ в $\mathbb{R}^n$, а все ее собственные значения вещественны и
    неотрицательны.
    Упорядочим их по невозрастанию:
    \[ \lambda_1 \ge \lambda_2 \ge \dots \ge \lambda_r > \lambda_{r+1} = \dots = \lambda_n = 0. \]
    При этом:
    \[ A^\top A v_i = \lambda_i v_i, \quad v_i^\top v_j = \delta_{ij}. \]
    Векторы $v_1, \dots, v_n$ образуют столбцы ортогональной матрицы $V \in \mathbb{R}^{n \times n}$ (
    $V^\top V = I_n$).
    \item \textbf{Определение сингулярных чисел:} Положим
    \[ \sigma_i \coloneqq \sqrt{\lambda_i} > 0 \quad \text{для } i = 1, \dots, r. \]
    \item \textbf{Построение левых сингулярных векторов $u_i$:} Для каждого $i \in \{1,\dots, r\}$ определим вектор:
    \[ u_i \coloneqq \frac{1}{\sigma_i} A v_i \in \mathbb{R}^m. \]
    Проверим скалярные произведения векторов $u_i$ и $u_j$:
    \[ u_i^\top u_j = \left(\frac{1}{\sigma_i} A v_i\right)^\top \left(\frac{1}{\sigma_j} A v_j\right) =
    \frac{1}{\sigma_i \sigma_j} v_i^\top (A^\top A v_j) = \frac{1}{\sigma_i \sigma_j} v_i^\top (\lambda_j v_j) =
    \frac{\lambda_j}{\sigma_i \sigma_j} (v_i^\top v_j). \]
    Так как $v_i^\top v_j = \delta_{ij}$ и $\lambda_i = \sigma_i^2$:
    \[ u_i^\top u_j = \frac{\sigma_j^2}{\sigma_i \sigma_j} \delta_{ij} = \delta_{ij}. \]
    Следовательно, система векторов $\{u_1, \dots, u_r\}$ строго ортонормирована в $\mathbb{R}^m$.
    \item Если $r < m$, дополним $\{u_1, \dots, u_r\}$ процедурой Грама-Шмидта до полного ортонормированного базиса
    $\{u_1, \dots, u_m\}$ пространства $\mathbb{R}^m$.
    \item Составим из них ортогональную матрицу $U = [u_1 \dots u_m]
    \in \mathbb{R}^{m \times m}$ ($U^\top U = I_m$).
    \item По построению:
    \[ A v_i = \sigma_i u_i \quad (1 \le i \le r), \quad A v_i = 0 \quad (r < i \le n). \]
    В матричной форме это записывается как:
    \[ A V = U \Sigma \iff A = U \Sigma V^\top, \] что полностью завершает конструктивный вывод SVD\@.
\end{enumerate}

% -------------------------------------------------------------------------

\subsection{Теоретический вопрос IV.2: Почему $A^\top A$ и $A A^\top$ имеют общие ненулевые eigenvalues?}
\textbf{Формулировка вопроса:} Докажите, что матрицы $A^\top A$ и $A A^\top$ имеют в точности одинаковые ненулевые
собственные значения с учётом кратностей.

\textbf{Ответ и строгое доказательство:}

Пусть $A \in \mathbb{R}^{m \times n}$.
Матрица $A^\top A$ имеет размер $n \times n$, а матрица $A A^\top$ имеет размер $m \times m$.
\begin{enumerate}
    \item Пусть $\lambda \ne 0$ — ненулевое собственное значение матрицы $A^\top A$, и $v \in \mathbb{R}^n$ —
    соответствующий ненулевой собственный вектор:
    \[ (A^\top A) v = \lambda v, \quad v \ne 0. \]
    \item Умножим обе части равенства слева на матрицу $A$:
    \[ A (A^\top A v) = A (\lambda v) \iff (A A^\top) (A v) = \lambda (A v). \]
    \item Определим вектор $w \coloneqq A v \in \mathbb{R}^m$.
    Проверим, отличен ли он от нуля:
    Вычислим квадрат евклидовой нормы вектора $w$:
    \[ \|w\|_2^2 = w^\top w = (A v)^\top (A v) = v^\top (A^\top A v) = v^\top (\lambda v) = \lambda \|v\|_2^2. \]
    Так как $\lambda > 0$ и $v \ne 0$, имеем $\|w\|_2^2 > 0$, откуда $w \ne 0$.
    \item Следовательно, вектор $w = Av$ является \textbf{ненулевым собственным вектором} матрицы $A A^\top$,
    отвечающим тому же самому собственному значению $\lambda$.
    \item Проводя аналогичное рассуждение в обратную сторону: если $(A A^\top) u = \mu u$ при $\mu \ne 0, u \ne 0$,
    то умножение слева на $A^\top$ даёт $(A^\top A)(A^\top u) = \mu (A^\top u)$ с ненулевым вектором $A^\top u$.
    \item \textbf{Совпадение геометрических кратностей:} Отображение $v \mapsto Av$ устанавливает линейный
    изоморфизм между собственным подпространством матрицы $A^\top A$ для значения $\lambda$ и собственным
    подпространством матрицы $A A^\top$ для того же $\lambda$.
    Следовательно, их размерности совпадают.
\end{enumerate}

% -------------------------------------------------------------------------

\subsection{Теоретический вопрос IV.3: Объясните outer-product SVD form}
\textbf{Формулировка вопроса:} Объясните представление SVD в виде суммы внешних произведений (outer-product form), ее
блочную алгебру и физический смысл.

\textbf{Ответ:}

Используя редуцированное сингулярное разложение $A = U_r \Sigma_r V_r^\top$, где $U_r = [u_1 \dots u_r]$, $V_r = [v_1
\dots v_r]$ и $\Sigma_r = \diag(\sigma_1, \dots, \sigma_r)$, распишем матричное произведение по правилу блочного
умножения «столбец на строку»:
\[ A = U_r \Sigma_r V_r^\top = [u_1 \dots u_r] \begin{pmatrix}
                                                           \sigma_1 v_1^\top \\ \vdots \\ \sigma_r v_r^\top
\end{pmatrix} = \sum_{i=1}^r \sigma_i u_i v_i^\top. \]

\textbf{Анализ элементарного слагаемого $u_i v_i^\top$:}
\begin{enumerate}
    \item $u_i \in \mathbb{R}^m$ ($m \times 1$), $v_i^\top \in \mathbb{R}^{1 \times n}$ ($1 \times n$).
    \item Произведение $u_i v_i^\top$ является матрицей размера $m \times n$ и имеет
    \textbf{ранг ровно 1} (столбцы пропорциональны $u_i$, строки пропорциональны $v_i^\top$).
    \item Матрицы $u_i v_i^\top$ взаимно ортогональны относительно скалярного произведения Фробениуса $\langle A, B
    \rangle_F = \tr(A^\top B)$:
    \[ \tr\left((u_i v_i^\top)^\top (u_j v_j^\top)\right) = \tr(v_i u_i^\top u_j v_j^\top) = (u_i^\top u_j)
    \tr(v_i v_j^\top) = \delta_{ij} (v_j^\top v_i) = \delta_{ij}. \]
\end{enumerate}

\textbf{Физическая и концептуальная интерпретация:} Внешняя форма SVD показывает, что любое сложное преобразование $A$
состоит из параллельной суммы $r$ элементарных независимых каналов:
\begin{itemize}
    \item $v_i^\top x$ — измерение проекции входа на входной базисный вектор $v_i$;
    \item $\sigma_i$ — коэффициент усиления (амплитуда, сила сигнала в данном канале);
    \item $u_i$ — синтез выходного направления в пространстве отклика.
\end{itemize}

\newpage

% =========================================================================
% ЧАСТЬ 2: ПРАКТИЧЕСКИЕ ВОПРОСЫ (PRACTICE)
% =====================================================================

    \chapter{Практические вопросы (Practice)}
\section{Модуль I. Векторные пространства и евклидова геометрия}

\subsection{Практический вопрос I.1: Как проверить, что subset является subspace?}
\textbf{Ответ:} Для проверки, является ли подмножество $U \subseteq V$ подпространством, применяется стандартный
трехшаговый алгоритм:
\begin{enumerate}
    \item \textbf{Проверка нуля:} Проверить, содержит ли $U$ нулевой вектор ($0 \in U$).
    Если $0 \notin U$, множество гарантированно не является подпространством.
    \item \textbf{Замкнутость по сложению:} Взять два произвольных элемента $u_1, u_2 \in U$ и проверить, что их
    сумма $u_1 + u_2 \in U$.
    \item \textbf{Замкнутость по умножению на скаляр:} Взять произвольный $u \in U$ и скаляр $c \in \mathbb{R}$
    и проверить, что $c u \in U$.
\end{enumerate}
\emph{Объединенный критерий:} $\forall u_1, u_2 \in U, \forall \alpha, \beta \in \mathbb{R} \implies \alpha u_1 + \beta
u_2 \in U$.

\subsection{Практический вопрос I.2: Как span, independence и basis работают в вычислениях?}
\textbf{Ответ:} В вычислительной практике все три концепции реализуются через метод Гаусса над матрицей $A$,
составленной из исследуемых векторов как из столбцов:
\begin{enumerate}
    \item \textbf{Проверка линейной независимости:} Приводим $A$ к ступенчатому виду.
    Если все столбцы содержат ведущие элементы (нет свободных переменных), система независима.
    \item \textbf{Нахождение базиса $\spanop(v_1, \dots, v_k)$:} Столбцы исходной матрицы $A$, соответствующие
    ведущим позициям ступенчатого вида, образуют базис оболочки.
    \item \textbf{Проверка принадлежности вектора $b \in \spanop(v_1, \dots, v_k)$:} Решаем систему $Ax = b$.
    Если система совместна , $b \in \spanop$; коэффициенты $x$ дают координаты разложения.
\end{enumerate}

\subsection{Практический вопрос I.3: Как меняются coordinates при смене basis?}
\textbf{Ответ:} Пусть вектор $x$ имеет координаты $[x]_B$ в базисе $B$ и $[x]_C$ в базисе $C$.
Строим матрицу перехода $P_{C \leftarrow B}$, столбцами которой являются координаты векторов старого базиса $B$,
выраженные в новом базисе $C$:
\[ [x]_C = P_{C \leftarrow B} [x]_B. \]
Обратный пересчет осуществляется с помощью обратной матрицы: $[x]_B = (P_{C \leftarrow B})^{-1} [x]_C$.

\subsection{Практический вопрос I.4: Как соотносятся dot product, norm, distance и angle?}
\textbf{Ответ:} Скалярное произведение $\langle x, y \rangle$ является первоисточником всех метрических величин:
\begin{itemize}
    \item \textbf{Норма (длина):} $\|x\| = \sqrt{\langle x, x \rangle}$;
    \item \textbf{Расстояние:} $d(x, y) = \|x - y\| = \sqrt{\langle x-y, x-y \rangle}$;
    \item \textbf{Угол:} $\cos \theta = \frac{\langle x, y \rangle}{\|x\| \|y\|}$ (корректность формулы
    $|\cos \theta| \le 1$ обеспечена неравенством Коши-Буняковского-Шварца);
    \item \textbf{Ортогональность:} $\theta = 90^\circ \iff \langle x, y \rangle = 0$.
\end{itemize}

\subsection{Практический вопрос I.5: Как работает Principal component analysis (PCA)?}
\textbf{Ответ:} Пошаговый алгоритм PCA:
\begin{enumerate}
    \item \textbf{Центрирование данных:} Вычесть среднее по каждому признаку: $X_c = X - \mu$.
    \item \textbf{Ковариационная матрица:} Вычислить $\Sigma=\frac{1}{n-1} X_c^\top X_c \in \mathbb{R}^{d \times d}$
    \item \textbf{Спектральное разложение / SVD:} Найти собственные векторы $v_i$ матрицы $\Sigma$
    (или правые сингулярные векторы $X_c$).
    \item \textbf{Сортировка:} Упорядочить компоненты по убыванию собственных значений $\lambda_1 \ge \lambda_2
    \dots \ge \lambda_d$.
    \item \textbf{Проекция:} Выбрать $k$ первых компонент $V_k = [v_1 \dots v_k]$ и вычислить сжатые координаты
    $Z = X_c V_k \in \mathbb{R}^{n \times k}$.
\end{enumerate}

\newpage

\section{Модуль II. Матрицы, линейные отображения, ранг и СЛУ}

\subsection{Практический вопрос II.1: Как построить matrix of a linear transformation?}
\textbf{Ответ:} Алгоритм построения матрицы отображения $T \colon V \to W$:
\begin{enumerate}
    \item Зафиксировать базис $\mathcal{B} = (v_1,\dots,v_n)$ в $V$ и базис $\mathcal{C} = (w_1, \dots, w_m)$ в $W$.
    \item Подействовать оператором $T$ на каждый базисный вектор $v_j$ ($j = 1, \dots, n$).
    \item Разложить полученный вектор $T(v_j)$ по базису $\mathcal{C}$: $T(v_j) = \sum_{i=1}^m a_{ij} w_i$.
    \item Записать найденные скаляры $a_{1j}, \dots, a_{mj}$ в качестве $j$-го столбца результирующей матрицы.
\end{enumerate}

\subsection{Практический вопрос II.2: Как determinant, rank и inverse диагностируют одну проблему?}
\textbf{Ответ:} Все три величины диагностируют проблему потери информации (вырожденности):
\begin{itemize}
    \item $\det A = 0 \iff$ объём сжимается в нуль;
    \item $\rank A < n \iff$ образ имеет размерность меньше размерности прообраза;
    \item Несуществование $A^{-1} \iff$ преобразование невозможно обратить вспять из-за схлопывания ядра в ноль.
\end{itemize}

\subsection{Практический вопрос II.3: Как определить число решений $Ax = b$?}
\textbf{Ответ:} По теореме Кронекера-Капелли приводим расширенную матрицу $(A \mid b)$ к ступенчатому виду:
\begin{enumerate}
    \item Если $\rank(A) < \rank(A \mid b)$ — \textbf{нет решений} (система несовместна, возникла строка вида
    $(0 \dots 0 \mid c)$ при $c \ne 0$).
    \item Если $\rank(A) = \rank(A \mid b) = n$ (числу неизвестных) — \textbf{ровно одно решение}.
    \item Если $\rank(A) = \rank(A \mid b) = r < n$ — \textbf{бесконечно много решений}
    (размерность аффинного пространства решений равна $n - r$).
\end{enumerate}

\subsection{Практический вопрос II.4: Почему LU полезно для systems?}
\textbf{Ответ:} Факторизация $A = LU$ требует $\mathcal{O}\left(\frac{2}{3} n^3\right)$ операций и выполняется один раз.
После этого решение системы $Ax = b$ для любой новой правой части $b$ сводится к решению двух треугольных систем ($Ly =
b$ и $Ux = y$), каждая из которых требует лишь $\mathcal{O}(n^2)$ операций.
Прямой метод Гаусса требовал бы $\mathcal{O}(n^3)$ на каждый вектор $b$.

\subsection{Практический вопрос II.5: Какие common dimensional mistakes возникают с matrices?}
\textbf{Ответ:}
\begin{enumerate}
    \item Умножение $A_{(m \times k)} B_{(p \times n)}$ при $k \ne p$.
    \item Сложение матриц разных размеров.
    \item Ошибочное транспонирование: $(AB)^\top \ne A^\top B^\top$ (верно $B^\top A^\top$).
    \item Путаница между внутренним произведением $x^\top y\in\mathbb{R}$ и внешним
    $x y^\top \in\mathbb{R}^{n \times n}$.
\end{enumerate}

\newpage

\section{Модуль III. Собственные значения и матричные разложения}

\subsection{Практический вопрос III.1: Как compute eigenvectors after eigenvalues?}
\textbf{Ответ:}
\begin{enumerate}
    \item Подставить найденное значение $\lambda_i$ в однородную систему $(A - \lambda_i I) v = 0$.
    \item Привести матрицу $A - \lambda_i I$ методом Гаусса к ступенчатому виду.
    \item Разделить переменные на базисные (ведущие) и свободные.
    \item Выразить базисные переменные через свободные и записать фундаментальную систему решений (ФСР), столбцы
    которой образуют базис собственного подпространства $E_{\lambda_i}$.
\end{enumerate}

\subsection{Практический вопрос III.2: Как диагонализация упрощает возведение в степень?}
\textbf{Ответ:} Если $A = P \Lambda P^{-1}$, то:
\[ A^k = (P \Lambda P^{-1})^k = P \Lambda^k P^{-1} = P \diag(\lambda_1^k, \dots, \lambda_n^k) P^{-1}. \]
Вместо $k - 1$ матричных умножений со сложностью $\mathcal{O}(k n^3)$ вычисляются лишь степени чисел $\lambda_i^k$, что
требует $\mathcal{O}(n^3 + n \log k)$ операций.

\subsection{Практический вопрос III.3: Как геометрически интерпретировать eigenvectors?}
\textbf{Ответ:} Собственный вектор $v$ указывает направление в пространстве, вдоль которого преобразование действует как
\textbf{чистое одномерное растяжение или сжатие} на коэффициент $\lambda$.
Вектор $Av$ остаётся на той же прямой, что и вектор $v$ (прямая инвариантна относительно оператора).

\subsection{Практический вопрос III.4: Почему symmetric matrices удобны?}
\textbf{Ответ:} 1) Все собственные значения гарантированно вещественны.
2) Собственные векторы ортогональны.
3) Всегда существует ортонормированный базис из собственных векторов ($A = Q \Lambda Q^\top$).
4) Обращение матрицы $Q$ тривиально и численно устойчиво: $Q^{-1} = Q^\top$.

\subsection{Практический вопрос III.5: Чем decompositions отличаются с точки зрения применимости?}
\textbf{Ответ:}
\begin{itemize}
    \item $LU$: общего назначения для квадратных матриц при решении СЛУ.
    \item Cholesky ($LL^\top$): в 2 раза быстрее $LU$, для симметричных положительно определенных матриц
    (МНК, оптимизация).
    \item $QR$: численно сверхустойчив, для переопределенных систем и вычисления собственных значений.
    \item Спектральное: для анализа квадратичных форм, физических мод и динамики.
    \item $SVD$: универсально для любых прямоугольных матриц, оптимально для низкоранговой аппроксимации и сжатия.
\end{itemize}

\newpage

\section{Модуль IV. SVD и итоговые инструменты}

\subsection{Практический вопрос IV.1: Опишите способ вычисления матрицы SVD}
\textbf{Ответ:}
\begin{enumerate}
    \item Вычислить симметричную матрицу $A^\top A$.
    \item Найти ее собственные значения $\lambda_i \ge 0$ и упорядочить $\lambda_1 \ge \dots \ge \lambda_r > 0$.
    \item Сингулярные числа: $\sigma_i = \sqrt{\lambda_i}$.
    \item Найти ортонормированные собственные векторы $v_i$ матрицы $A^\top A$ $\implies$ матрица $V$.
    \item Вычислить левые сингулярные векторы: $u_i = \frac{1}{\sigma_i} A v_i$ для $i = 1, \dots, r$.
    \item При необходимости дополнить $\{u_i\}$ до ортонормированного базиса $\mathbb{R}^m \implies$ матрица $U$.
    \item Собрать $A = U \Sigma V^\top$.
\end{enumerate}

\subsection{Практический вопрос IV.2: Как reduced SVD используется в рекомендательных системах?}
\textbf{Ответ:} Матрица оценок пользователей $R_{m \times n}$ факторизуется через усеченное SVD: $R \approx U_k \Sigma_k
V_k^\top = P Q^\top$.
Строки матрицы $P = U_k \Sigma_k^{1/2}$ отражают вкусовые профили пользователей в $k$ латентных темах, а строки $Q = V_k
\Sigma_k^{1/2}$ — профили контента.
Пропущенная оценка восстанавливается как $\hat{r}_{ui} = \langle p_u, q_i \rangle$.

\subsection{Практический вопрос IV.3: Чем полезна low-rank approximation? Назовите пример}
\textbf{Ответ:} Низкоранговая аппроксимация сохраняет главную энергию матрицы, отсекая шум и радикально сокращая память
с $m \cdot n$ до $k(m + n + 1)$.
\emph{Пример:} Сжатие изображений: фотография $1000 \times 1000$ при $k = 50$ занимает в 10 раз меньше памяти при почти
неотличимом качестве.

\subsection{Практический вопрос IV.4: Что такое pseudoinverse matrix, чем она полезна?}
\textbf{Ответ:} Псевдообратная матрица Мура-Пенроуза $A^+ = V \Sigma^+ U^\top$ обобщает понятие обратной матрицы на
случай любых прямоугольных или вырожденных матриц.
Она даёт наилучшее обобщенное решение системы $Ax = b$: $x^* = A^+ b$.
Если решений нет, $x^*$ минимизирует ошибку $\|Ax - b\|_2$; если решений бесконечно много, $x^*$ выбирает решение с
наименьшей нормой $\|x\|_2$.

\newpage

% =========================================================================
% ЧАСТЬ 3: БАЗОВЫЕ ВОПРОСЫ (BASIC)
% =========================================================================

    \chapter{Базовые вопросы (Basic)}
\section{Модуль I. Векторные пространства и евклидова геометрия}

\subsection{Базовый вопрос I.1: Дайте определение вещественного векторного пространства}
\textbf{Ответ:} Множество $V$ с операцией сложения $+ \colon V \times V \to V$ и умножения на скаляр $\cdot \colon
\mathbb{R} \times V \to V$, удовлетворяющее 8 аксиомам: 1) $u+v = v+u$; 2) $(u+v)+w = u+(v+w)$; 3) $\exists 0 \colon
v+0=v$; 4) $\exists (-v) \colon v+(-v)=0$; 5) $1 \cdot v = v$; 6) $\alpha(\beta v) = (\alpha\beta)v$; 7) $\alpha(u+v) =
\alpha u + \alpha v$; 8) $(\alpha+\beta)v = \alpha v + \beta v$.

\subsection{Базовый вопрос I.2: Что такое подпространство?}
\textbf{Ответ:} Непустое подмножество $U \subseteq V$, которое само является векторным пространством относительно тех же
операций.
Критерий: $0 \in U$ и $\forall u, v \in U, \alpha, \beta \in \mathbb{R} \implies \alpha u + \beta v \in U$.

\subsection{Базовый вопрос I.3: Что такое linear combination, span и generating set?}
\textbf{Ответ:} Линейная комбинация — выражение вида $\sum_{i=1}^k c_i v_i$.
Линейная оболочка $\spanop(S)$ — множество всех линейных комбинаций векторов из $S$.
Порождающая система — множество векторов $S$, для которого $\spanop(S) = V$.

\subsection{Базовый вопрос I.4: Определения linear independence, basis, dimension, coordinates}
\textbf{Ответ:}
\begin{itemize}
    \item \textbf{Линейная независимость:} $\sum c_i v_i = 0 \iff \text{все } c_i = 0$.
    \item \textbf{Базис:} линейно независимая порождающая система.
    \item \textbf{Размерность $\dim V$:} число векторов в базисе.
    \item \textbf{Координаты:} единственный набор скаляров $(x_1, \dots, x_n)$ такой, что $x = \sum x_i e_i$.
\end{itemize}

\subsection{Базовый вопрос I.5: Что такое inner product, norm, distance и orthogonality?}
\textbf{Ответ:} Скалярное произведение — симметричная, положительно определенная билинейная форма.
Норма — длина вектора $\|x\| = \sqrt{\langle x, x \rangle}$.
Расстояние — метрика $d(x, y) = \|x - y\|$.
Ортогональность — условие $\langle x, y \rangle = 0$.

\subsection{Базовый вопрос I.6: Что такое orthogonal projection?}
\textbf{Ответ:} Ортогональная проекция вектора $v$ на подпространство $U$ — это единственный вектор $p \in U$, такой что
вектор ошибки $v - p$ ортогонален подпространству $U$: $(v - p) \perp U$.

\subsection{Базовый вопрос I.7: Что такое hyperplane и normal vector?}
\textbf{Ответ:} Гиперплоскость в $\mathbb{R}^n$ — аффинное подпространство размерности $n - 1$, задаваемое уравнением
$\langle n, x \rangle = c$.
Вектор $n \ne 0$ называется вектором нормали (он перпендикулярен гиперплоскости).

\newpage

\section{Модуль II. Матрицы, линейные отображения, ранг и СЛУ}

\subsection{Базовый вопрос II.1: Что такое matrix multiplication?}
\textbf{Ответ:} Операция $C = AB$, где $c_{ij} = \sum_{k=1}^p a_{ik} b_{kj}$ (скалярное произведение $i$-й строки $A$ на
$j$-й столбец $B$).
Является координатной записью композиции операторов.

\subsection{Базовый вопрос II.2: Что такое linear transformation и как matrix его представляет?}
\textbf{Ответ:} Отображение $T \colon V \to W$, сохраняющее линейные комбинации: $T(\alpha u + \beta v) = \alpha T(u) +
\beta T(v)$.
Столбцы матрицы оператора представляют собой координаты образов базисных векторов.

\subsection{Базовый вопрос II.3: Определения inverse, determinant, rank и nullspace}
\textbf{Ответ:}
\begin{itemize}
    \item \textbf{Обратная матрица $A^{-1}$:} удовлетворяет $A A^{-1} = A^{-1} A = I$.
    \item \textbf{Определитель $\det A$:} знакопеременная полилинейная функция строк, коэффициент изменения
    ориентированного объёма.
    \item \textbf{Ранг $\rank A$:} число линейно независимых столбцов (размерность образа).
    \item \textbf{Ядро $\Null A$:} множество решений уравнения $Ax = 0$.
\end{itemize}

\subsection{Базовый вопрос II.4: Сформулируйте rank-nullity theorem}
\textbf{Ответ:} Для линейного отображения $T \colon V \to W$:
\[ \dim V = \rank(T) + \nullity(T). \]

\subsection{Базовый вопрос II.5: Как записывается и интерпретируется СЛУ?}
\textbf{Ответ:} Запись $Ax = b$.
Интерпретация строк (row picture): пересечение $m$ гиперплоскостей.
Интерпретация столбцов (column picture): представление вектора $b$ в виде линейной комбинации столбцов матрицы $A$.

\subsection{Базовый вопрос II.6: Что такое Gaussian elimination?}
\textbf{Ответ:} Алгоритм последовательного исключения неизвестных с помощью элементарных преобразований строк для
приведения матрицы к верхнетреугольному (ступенчатому) виду.

\subsection{Базовый вопрос II.7: Что такое LU decomposition?}
\textbf{Ответ:} Факторизация матрицы в виде $A = LU$, где $L$ — нижнетреугольная матрица с единицами на диагонали, а $U$
— верхнетреугольная матрица.

\newpage

\section{Модуль III. Собственные значения и матричные разложения}

\subsection{Базовый вопрос III.1: Дайте определение eigenvalue и eigenvector}
\textbf{Ответ:} Ненулевой вектор $v \ne 0$ и скаляр $\lambda$ такие, что $Av = \lambda v$.

\subsection{Базовый вопрос III.2: Как найти eigenvalues и eigenvectors?}
\textbf{Ответ:} 1) Найти корни характеристического уравнения $\det(A - \lambda I) = 0$.
2) Для каждого $\lambda$ найти ненулевые решения системы $(A - \lambda I)v = 0$.

\subsection{Базовый вопрос III.3: Что такое algebraic и geometric multiplicity?}
\textbf{Ответ:} Алгебраическая кратность — кратность корня в многочлене $\det(A - \lambda I) = 0$.
Геометрическая кратность — размерность ядра $\dim \ker(A - \lambda I)$.

\subsection{Базовый вопрос III.4: Что такое eigendecomposition и diagonalization?}
\textbf{Ответ:} Представление $A=P \Lambda P^{-1}$, где столбцы $P$ — собственные векторы, а диагональ $\Lambda$ —
собственные значения.
Возможно тогда и только тогда, когда существует базис из собственных векторов.

\subsection{Базовый вопрос III.5: Сформулируйте real spectral theorem}
\textbf{Ответ:} Всякая вещественная симметричная матрица $A = A^\top$ имеет вещественный спектр и ортогонально
диагонализируема: $A = Q \Lambda Q^\top$, где $Q$ — ортогональная матрица ($Q^\top Q = I$).

\subsection{Базовый вопрос III.6: Что такое orthogonal matrix?}
\textbf{Ответ:} Квадратная матрица $Q$, для которой $Q^\top Q = Q Q^\top = I$ ($Q^{-1} = Q^\top$).
Сохраняет длины векторов и углы ($\|Qx\| = \|x\|$).

\subsection{Базовый вопрос III.7: Что такое positive definite, QR и Cholesky?}
\textbf{Ответ:}
\begin{itemize}
    \item \textbf{Positive definite:} $x^\top A x > 0$ для всех $x \ne 0$ (все $\lambda_i > 0$).
    \item \textbf{QR-разложение:} $A = QR$ ($Q$ ортонормирована, $R$ верхнетреугольная).
    \item \textbf{Разложение Cholesky:} $A = L L^\top$ ($L$ нижнетреугольная с положительной диагональю).
\end{itemize}

\newpage

\section{Модуль IV. SVD и итоговые инструменты}

\subsection{Базовый вопрос IV.1: Что такое SVD?}
\textbf{Ответ:} Сингулярное разложение произвольной матрицы: $A = U \Sigma V^\top$, где $U$ и $V$ — ортогональные
матрицы, а $\Sigma$ содержит неотрицательные сингулярные числа на диагонали.

\subsection{Базовый вопрос IV.2: Как singular values связаны с eigenvalues?}
\textbf{Ответ:} Сингулярные числа матрицы $A$ равны квадратным корням из собственных значений матрицы $A^\top A$:
$\sigma_i(A) = \sqrt{\lambda_i(A^\top A)}$.
Для положительно полуопределенной матрицы сингулярные числа совпадают с собственными значениями.

\subsection{Базовый вопрос IV.3: Что такое reduced SVD?}
\textbf{Ответ:} Компактная форма $A = U_r \Sigma_r V_r^\top$, где отброшены нулевые сингулярные числа и соответствующие
столбцы матриц $U$ и $V$.
Размерность $\Sigma_r$ равна $r \times r$, где $r = \rank A$.

\subsection{Базовый вопрос IV.4: Каков geometric meaning of SVD?}
\textbf{Ответ:} Отображение единичной сферы в эллипсоид в три этапа: поворот в пространстве прообразов ($V^\top$),
масштабирование по осям на величины $\sigma_i$ ($\Sigma$), и поворот полученного эллипсоида в пространстве образов
($U$).

\subsection{Базовый вопрос IV.5: Что такое low-rank approximation и pseudoinverse?}
\textbf{Ответ:}
\begin{itemize}
    \item \textbf{Low-rank approximation:} усеченная сумма $A_k = \sum_{i=1}^k \sigma_i u_i v_i^\top$ ($k < r$),
    наилучшее приближение матрицы $A$ рангом $k$.
    \item \textbf{Pseudoinverse (Мура-Пенроуза):} матрица $A^+ = V \Sigma^+ U^\top$, обеспечивающая обобщенное
    решение задачи наименьших квадратов $x = A^+ b$ с минимальной евклидовой нормой.
\end{itemize}

\end{document}
